% Generated draft: requires mathematical and source review.
\documentclass[11pt]{article}
\usepackage[margin=25mm]{geometry}
\usepackage{fontspec}
\setmainfont{Latin Modern Roman}
\usepackage{amsmath,amssymb,mathtools,unicode-math}
\setmathfont{Latin Modern Math}
\usepackage{enumitem,needspace,xurl,hyperref}
\setlength{\parindent}{0pt}
\setlength{\parskip}{5pt}
\setlength{\emergencystretch}{3em}
\newcommand{\N}{\mathbb{N}}
\newcommand{\Z}{\mathbb{Z}}
\newcommand{\Q}{\mathbb{Q}}
\newcommand{\R}{\mathbb{R}}
\newcommand{\C}{\mathbb{C}}
\newcommand{\F}{\mathbb{F}}
\newcommand{\A}{\mathbb{A}}
\newcommand{\PP}{\mathbb{P}}
\newcommand{\E}{\mathbb{E}}
\newcommand{\eps}{\varepsilon}
\newcommand{\dd}{\,\mathrm{d}}
\begin{document}
\section*{UnsolvedMath ID: 700937}
\textbf{Source problem code:} catalogue entry
\subsection{Definitions and mathematical statement}
Let \(n\geq 1\) be an integer, let \(G_1,\ldots,G_n\) be groups, and write \(\pi_i:G_1\times\cdots\times G_n\to G_i\) for the coordinate projections. Suppose that
\[
H\trianglelefteq Q\leq G_1\times\cdots\times G_n,
\qquad \pi_i(H)=\pi_i(Q)=G_i\quad(1\leq i\leq n).
\]
Thus both \(H\) and \(Q\) are subdirect subgroups of the same finite direct product. Put \(P=Q/H\), and let \(p\) be a prime.

A finite \(p\)-group is a finite group of order \(p^m\) for some integer \(m\geq 0\). Such a group \(P\) is called regular if, for every \(a,b\in P\), there are an integer \(r\geq 0\) and elements \(c_1,\ldots,c_r\in [K,K]\), where \(K=\langle a,b\rangle\) and \([K,K]\) is its commutator subgroup, such that
\[
(ab)^p=a^p b^p c_1^p\cdots c_r^p.
\]
The product is the identity when \(r=0\).

The supplied question is: if \(P=Q/H\) is a \(p\)-group, must it be a regular \(p\)-group? Under the conventional finite interpretation, this asks whether, for every prime \(p\), every integer \(n\geq 1\), and every choice of groups and subgroups satisfying the displayed hypotheses, a quotient \(Q/H\) of finite \(p\)-power order satisfies the regularity condition above. No finiteness assumption on the factors, \(Q\), or \(H\) is stated.

The source does not explicitly require \(Q/H\) to be finite or define regularity. The term \(p\)-group can also allow infinite groups in which every element has finite \(p\)-power order. Consequently, the finite formulation above is a conditional interpretation; an intended infinite-group formulation and its regularity convention remain unresolved.
\subsection{Short English statement}
If a group and a normal subgroup both project onto every factor of a finite direct product, must their quotient, when a \(p\)-group, be regular? The supplied record leaves the finiteness and regularity conventions unspecified.
\subsection{Sources}
\begin{itemize}
\item Local UnsolvedMath registry, permanent ID \(700937\), exact source problem code \texttt{catalogue entry}; supplied complete source record, upstream revision \texttt{d371ee917ac9647fbb06cd6fc529deec6f3dda82}.
\item The Kourovka Notebook (Unsolved Problems in Group Theory), original Problem 13.19, 13th issue (1995); proposed by Yu. M. Gorchakov.
\item The supplied background also identifies the 21st issue edition (2026), \texttt{arXiv:1401.0300v48}. The registry collection identifies E. I. Khukhro and V. D. Mazurov as editors of the notebook covering Issues 1--21 (1965--2026).
\item The registry labels this problem open and records an update dated 2026-10-16. This is supplied source metadata, not an independently verified assessment of the mathematical or literature status.
\end{itemize}
\end{document}
